\emph{Scope.} Clauses (i), (ii) and (iv) use only the mode functions $\chi_{k}$ of a homogeneous Hadamard reference state $\omega_{\mathrm{ref}}$ (smooth in $\eta$, continuous in $k > 0$) and the smooth-difference characterisation of Fact \claimref{synthesis.imported-microlocal-hadamard-condition}: the wave front set condition of part (a), with part (b)'s statement that two Hadamard states differ by a smooth kernel. Below they are written for the checked instance, the conformal vacuum $\omega_{0}$ on the frozen point's field, with $e^{-2ik\eta_{d}}$ in place of the general phase $\theta_{k} = \chi_{k}(\eta_{d})^{2}/|\chi_{k}(\eta_{d})|^{2}$; for a general reference replace $e^{-ik\eta}/\sqrt{2k}$ by $\chi_{k}$ throughout, and the bound on the integrand in (i) by the supremum of $|\chi_{k}|$ and its $\eta$-derivatives over the compact shell and a compact time interval. Clause (iii) uses the conformal modes; on a spatially flat Robertson-Walker field other than the frozen point's, the conformal vacuum of the conformally coupled field is the flat vacuum of $\chi = a\phi$ by the conformal covariance of Fact \claimref{synthesis.imported-conformal-covariance-of-massless-fields-and-local-conformal-flatness}(b), and that it is Hadamard there is the assumption stated in (iii). These ingredients are standard (finite-band Bogoliubov data, rapid decay of the occupations of a Hadamard state, Ford and Parker's infrared divergence, Fact \claimref{synthesis.imported-ford-parker-infrared-divergence}, the alpha-vacua); the scalar necessity argument of (iii) is the counterpart of the occupation-moment criterion for the massless Dirac field (\claimref{synthesis.deposit-over-the-conformal-vacuum-is-hadamard-exactly-when-every-occupation-moment-is-finite}).

\emph{Setting.} On the frozen point's field in its flat slicing, $a = -1/(H\eta)$ and $a''/a = 2/\eta^{2}$. A scalar of weight $w$ in the dictionary of Fact \claimref{standard-model.imported-frozen-point-dictionary} has $m^{2} = -w(w + 3)H^{2}$, and $\chi = a\phi$ obeys $\chi'' + (k^{2} + m^{2}a^{2} - a''/a)\chi = 0$. For $w = -1$, $m^{2}a^{2} = 2H^{2}/(H^{2}\eta^{2}) = 2/\eta^{2}$ cancels $a''/a$ exactly, so $\chi'' + k^{2}\chi = 0$ (for the massless point $w = 0$ the term $-2/\eta^{2}$ stays; the check verifies both). The solutions $\chi_{k} = e^{-ik\eta}/\sqrt{2k}$ have Wronskian $\chi_{k}\bar\chi_{k}' - \chi_{k}'\bar\chi_{k} = i$. The state with these modes is invariant under the stabiliser of the frozen point and is Hadamard; by Fact \claimref{synthesis.imported-de-sitter-invariant-hadamard-state-unique} ($m^{2} = 2H^{2} > 0$) it is the unique such state, the one used in Theorem \claimref{synthesis.frozen-point-geometry-reads-descent-as-vacuum}(iii), whose slice spectrum is $\Delta_{0}^{2} = k^{3}|\chi_{k}|^{2}/(2\pi^{2}a^{2}) = k^{2}/(4\pi^{2}a^{2})$. The complex unit of the mode functions is read as $I_{4}$, $I_{4}^{2} = -1$; the check verifies that the product of numbers $a + bI_{4}$ in the spacetime algebra is the complex product it uses.

A homogeneous isotropic quasi-free state $\omega$ has, in the Fock modes $a_{\vec k}$ of $\omega_{0}$, $\omega(a_{\vec k}^{*}a_{\vec k'}) = n_{k}\delta(\vec k - \vec k')$ and $\omega(a_{\vec k}a_{-\vec k'}) = c_{k}\delta(\vec k - \vec k')$; positivity of the two-mode covariance requires $n_{k} \ge 0$ and $|c_{k}|^{2} \le n_{k}(n_{k} + 1)$, with equality exactly for pure states, where $n_{k} = |\beta_{k}|^{2}$ and $c_{k} = \alpha_{k}\bar\beta_{k}$ up to the phase convention. The symmetrised equal-time two-point function of $\chi$ at $\eta_{d}$ then has Fourier weight $|\chi_{k}|^{2}(1 + 2n_{k}) + 2\,\mathrm{Re}(c_{k}\chi_{k}(\eta_{d})^{2})$, which is $|\chi_{k}|^{2}\rho_{k}$ with $\rho_{k} = 1 + 2n_{k} + 2\,\mathrm{Re}(c_{k}e^{-2ik\eta_{d}})$. For a pure state this is $|u_{k}|^{2} = |\chi_{k}|^{2}|\alpha_{k} + \beta_{k}e^{2ik\eta_{d}}|^{2}$.

\emph{(i)} Write the two-point function of $\omega$ minus that of $\omega_{0}$ as the mode integral
\[
D(x, y) = \frac{1}{a(\eta)a(\eta')}\int_{k_{1} \le |\vec k| \le k_{2}}\frac{d^{3}k}{(2\pi)^{3}}\Bigl[2n_{k}\,\mathrm{Re}\bigl(\chi_{k}(\eta)\bar\chi_{k}(\eta')\bigr)\cos(\vec k\cdot(\vec x - \vec y)) + \bigl(c_{k}\chi_{k}(\eta)\chi_{k}(\eta') + \text{c.c.}\bigr)\cos(\vec k\cdot(\vec x - \vec y))\Bigr]
\]
(the antisymmetric part, the commutator, is state-independent). The integrand is bounded by $(2\sup n + 2\sup|c|)/(2k_{1})$ on a compact set of $\vec k$, it is smooth in $(\eta, \eta', \vec x, \vec y)$, and each derivative of order $N$ multiplies it by at most $k_{2}^{N}$; so differentiation under the integral is allowed to every order and $D$ is $C^{\infty}$ (indeed real-analytic). It is a bisolution, being a superposition of products of solutions. By Fact \claimref{synthesis.imported-microlocal-hadamard-condition}(a) the Hadamard property is the wave front set condition on the two-point distribution, and adding a smooth kernel does not change a wave front set; $\omega_{0}$ satisfies the condition, hence so does $\omega$. The check computes an instance: for $n_{k} = n$ constant and $c_{k} = 0$ on $[k_{1}, k_{2}]$ the equal-time difference is $n(\cos k_{1}r - \cos k_{2}r)/(2\pi^{2}a^{2}r^{2})$, from $\int_{k_{1}}^{k_{2}}k\,\sin(kr)/(kr)\,dk = (\cos k_{1}r - \cos k_{2}r)/r^{2}$, whose Taylor series in $r^{2}$ it verifies term by term against the term-by-term integral to order $r^{22}$; its value at $r = 0$ is $n(k_{2}^{2} - k_{1}^{2})/(4\pi^{2}a^{2})$, while $\omega_{0}$'s equal-time kernel is $1/(4\pi^{2}a^{2}r^{2})$.

For the realisation of a ratio, put $q_{k} = \sqrt{\rho_{k}}$, $\alpha_{k} = (q_{k} + 1/q_{k})/2 = \cosh s_{k}$ and $\beta_{k} = \pm\frac12(q_{k} - 1/q_{k})e^{-2ik\eta_{d}}$, so $|\beta_{k}| = \sinh|s_{k}|$ with $s_{k} = \ln q_{k}$ and $|\alpha_{k}|^{2} - |\beta_{k}|^{2} = 1$. With the upper sign $|\alpha_{k} + \beta_{k}e^{2ik\eta_{d}}|^{2} = (\cosh s_{k} + \sinh s_{k})^{2} = q_{k}^{2} = \rho_{k}$, with the lower sign $(\cosh s_{k} - \sinh s_{k})^{2} = 1/q_{k}^{2}$; choosing the sign by whether $\rho_{k} \ge 1$ (and using $|\sinh|$) gives the stated state. The coefficients are bounded on the band because $\rho$ is continuous and positive on a compact interval, and set to zero off it, so the first part applies. The check verifies, with exact Gaussian rationals and $2k$ a rational square, that for nine values of $q$ (from $1/3$ to $1000$), both signs, four values of $k$ and six points $e^{-ik\eta_{d}}$ of the unit circle, the new modes keep the Wronskian $i$, $|\beta|^{2} = ((q - 1/q)/2)^{2}$, and $|u_{k}|^{2}/|\chi_{k}|^{2}$ equals $q^{2}$ or $1/q^{2}$ exactly. Since the ratio is arbitrary, so are the amplitude and the local index $d\ln\Delta_{\omega}^{2}/d\ln k = 2 + d\ln\rho_{k}/d\ln k$ on the band.

\emph{(ii)} Write $t = 2\,\mathrm{Re}(c_{k}e^{-2ik\eta_{d}})$, so $\rho = 1 + 2n + t$ and $t^{2} \le 4|c|^{2} \le 4n(n + 1)$. Then $(\rho - 1)^{2} - 4n\rho = (2n + t)^{2} - 4n(1 + 2n + t) = t^{2} - 4n^{2} - 4n \le 0$. Equality holds when $t^{2} = 4n(n + 1)$, which the states of (i) attain ($n = \sinh^{2}s$, $|c| = \sinh|s|\cosh s$, phase aligned or opposed). Solving the quadratic, $4n \ge \rho + 1/\rho - 2 = (\sqrt\rho - 1/\sqrt\rho)^{2}$. The check tests the inequality at every admissible sample of $n$ (five values), $c$ (twenty values, those with $|c|^{2} \le n(n + 1)$ kept) and six phases, and the equality for the pure data.

\emph{(iii)} Let $\omega$ be Hadamard. By part (b) of the Fact (the two-point functions of two Hadamard states differ by a smooth kernel; for the Klein-Gordon field this follows from part (a), both two-point distributions having the wave front set $C^{+}$ and the same antisymmetric part) $D$ is smooth. Apply $\partial_{\eta}\partial_{\eta'} + \nabla_{\vec x}\cdot\nabla_{\vec y}$ to $a(\eta)a(\eta')D$ and set $\eta = \eta'$: the result $e(\vec x - \vec y)$ is a smooth function of the separation. Its Fourier weight is $(|u_{k}'|^{2} + k^{2}|u_{k}|^{2}) - (|\chi_{k}'|^{2} + k^{2}|\chi_{k}|^{2})$ for pure states and, in general, $2n_{k}(|\chi_{k}'|^{2} + k^{2}|\chi_{k}|^{2}) + 2\,\mathrm{Re}\bigl(c_{k}(\chi_{k}'^{2} + k^{2}\chi_{k}^{2})\bigr)$. For the conformal modes $\chi_{k}'^{2} + k^{2}\chi_{k}^{2} = (-k^{2} + k^{2})\chi_{k}^{2} = 0$ exactly and $|\chi_{k}'|^{2} + k^{2}|\chi_{k}|^{2} = k$, so the weight is $2kn_{k} \ge 0$ with no cross term (the check verifies $|u'|^{2} + k^{2}|u|^{2} = k(|\alpha|^{2} + |\beta|^{2})$ for three pairs at every sample, and, as a control, that the massless point's Bunch-Davies mode $e^{-iz}(1 - i/z)$ has $\chi'^{2} + k^{2}\chi^{2} = \tfrac{k}{2}(2/z^{2} - 2i/z^{3} - 1/z^{4})e^{-2iz} \neq 0$). A function continuous at $0$ whose Fourier transform is a nonnegative function has that transform integrable, with integral the value at $0$ (smear with a normalised Gaussian of width $\epsilon$ and let $\epsilon \to 0$: the left side tends to $e(0)$, the right side increases to $\int 2kn_{k}\,d^{3}k/(2\pi)^{3}$ by monotone convergence). So $\int k^{3}n_{k}\,dk < \infty$. Applying the Laplacian $j$ times before setting the arguments equal multiplies the weight by $k^{2j}$ and gives $\int k^{3 + 2j}n_{k}\,dk < \infty$.

If $\rho_{k} = Ck^{q}$ for $k \ge K$, (ii) gives $4n_{k} \ge Ck^{q} + C^{-1}k^{-q} - 2$, which grows like $k^{|q|}$ for $q \neq 0$ and is the positive constant $C + C^{-1} - 2$ for $q = 0$, $C \neq 1$; in both cases $\int^{\infty}k^{3}n_{k}\,dk$ diverges, so $\omega$ is not Hadamard. The check computes exact lower bounds of $\int k^{3}n_{k}\,dk$ over the shells $[b^{j}, b^{j+1}]$, $b = 2^{10000}$, $j = 1, 2, 3$, for $q = -20351/10000$ (the law $n_{k} \propto k^{n_{s} - 3}$ at the stored central $n_{s} = 0.9649$, used only as a sample exponent), $q = -1$ and $q = 1/2$, and finds each above $1$ and increasing. A slice correlator homogeneous of degree $-2\Delta$ has $\Delta^{2} \propto k^{2\Delta}$ at every $k$ (the scaling lemma of Theorem \claimref{synthesis.frozen-point-geometry-reads-descent-as-vacuum}(iv)), so $\rho_{k} \propto k^{2\Delta - 2}$, excluded for $\Delta \neq 1$. The frozen point's dilatation $(\eta, \vec x) \mapsto (\lambda\eta, \lambda\vec x)$ fixes $\omega_{0}$ and acts on the modes by $\vec k \mapsto \vec k/\lambda$ with $k\eta$ fixed, so a dilatation-invariant $\omega$ has $n_{k}$ and $c_{k}$ constant on $k > 0$; then $\int k^{3}n\,dk < \infty$ forces $n = 0$, hence $c = 0$ and $\omega = \omega_{0}$ (the check shows the shell bounds of a constant occupation grow by more than a factor $15$ per shell of ratio $16$).

\emph{(iv)} For a test function $g$ on the slice, $\omega(\phi(g)^{2}) = \int\frac{d^{3}k}{(2\pi)^{3}}|\hat g(\vec k)|^{2}\frac{2\pi^{2}}{k^{3}}\Delta_{\omega}^{2}(k) = \int_{0}^{\infty}\frac{dk}{k}\Delta_{\omega}^{2}(k)\,\langle|\hat g|^{2}\rangle_{k}$, with $\langle|\hat g|^{2}\rangle_{k}$ the angular mean. Since $\hat g$ is continuous with $\hat g(0) = \int g$, for every $g$ of nonzero mean $\langle|\hat g|^{2}\rangle_{k} \ge \tfrac12|\int g|^{2}$ on some $(0, \delta)$, so finiteness for all $g$ needs $\int_{0}\Delta_{\omega}^{2}\,dk/k < \infty$. For $\Delta_{\omega}^{2} = Ck^{n - 1}$ on $(0, k_{1})$ the integral is $C\int_{0}^{k_{1}}k^{n - 2}\,dk$, divergent for $n \le 1$; for $\omega_{0}$, $\int_{0}k\,dk$ converges. The condition excludes only $n \le 1$ near $k = 0$: for instance $\rho_{k} = 1/k$ on $(0, 1]$ (slice index $2$) is realised by the pure states of (i) with $n_{k} = (1 - k)^{2}/(4k)$, $|c_{k}| = (1 - k^{2})/(4k)$; then $\int_{0}^{1}kn_{k}\,dk = 1/12$ and $\int_{0}^{1}k|c_{k}|\,dk = 1/6$, so the difference kernel is smooth (derivatives add bounded powers of $k$ on a bounded set), the state is Hadamard, and $\int_{0}\Delta_{\omega}^{2}\,dk/k$ converges: a power law on a band reaching $k = 0$ is allowed when it is blue enough, so the infrared clause concerns red and scale-invariant laws only. The check evaluates $\int_{\epsilon}^{1}k^{s - 1}\,dk = (1 - \epsilon^{s})/s$ exactly at $\epsilon = 2^{-10000m}$ for $s = -351/10000$ (values $(2^{351m} - 1)\cdot 10000/351$, unbounded), sums forty dyadic shells for $s = 0$ (total $20$, each $\tfrac12$), and bounds the vacuum's $s = 2$ by $\tfrac12$. The Hadamard condition (i) and (iii) involves only the behaviour of $n_{k}$ and $c_{k}$ as $k \to \infty$, so it does not decide this integral.

\emph{Conclusion.} By (i) any bounded band can carry any positive spectrum in a Hadamard state; by (iii) a red (or any power-law) ratio has to stop at some finite $k_{\mathrm{UV}}$; by (iv) a red law has to stop at some $k_{\mathrm{IR}} > 0$ for the state to exist on the slice fields of nonzero mean, while a blue law may continue to $k = 0$. Neither (i) nor (iii) nor (iv) gives a value of $k_{\mathrm{UV}}$, $k_{\mathrm{IR}}$, the amplitude or the index on the band.
